\chapter{The Curry--Howard--Lambek correspondence}

From this chapter onward, we will only consider cartesian closed categories, for which we will use the abbreviations CCC and CCCs.

\section{Natural number objects in CCCs}
The aim of this section is to provide cartesian closed categories with a notion of recursion,
which is fundamental if we want to interpret a simply $\lambda$-calculus with iterator in a CCC.
Even though this notion won't be necessary for the correspondence used to prove Lawvere's fixed-point theorem,
it is stil necessary to prove the Curry--Howard--Lambek correspondence in its general version, and nevertheless
interesting in its own right. One notable application of this notion is, for example, the interpretation of all primitive recursive functions in a CCC with a natural number object, see \cite[Part~III, \S\S~1--2]{LambekScott1994}.

\begin{definition}[Natural number object]
    Given a cartesian closed category $\mathcal{C}$, a \emph{natural number object} (NNO) is an initial object
    \[ 1 \xrightarrow{0} N \xrightarrow{s} N \]
    in the category of diagrams of the form $1 \xrightarrow{a} A \xrightarrow{b} A$ in $\mathcal{C}$. An equivalent way to define it is
    by saying that, for all such diagrams, there exists a unique arrow $h : N \to A$ such that the following diagram commutes:
    % https://q.uiver.app/#q=WzAsNixbMCwwLCIxIl0sWzEsMCwiTiJdLFsyLDAsIk4iXSxbMSwxLCJBIl0sWzIsMSwiQSJdLFswLDEsIjEiXSxbMCwxLCIwIl0sWzEsMiwicyJdLFszLDQsImYiXSxbNSwzLCJhIl0sWzAsNSwiIiwxLHsib2Zmc2V0IjotMSwic3R5bGUiOnsiaGVhZCI6eyJuYW1lIjoibm9uZSJ9fX1dLFsxLDMsImgiLDIseyJzdHlsZSI6eyJib2R5Ijp7Im5hbWUiOiJkYXNoZWQifX19XSxbMiw0LCJoIiwyLHsic3R5bGUiOnsiYm9keSI6eyJuYW1lIjoiZGFzaGVkIn19fV0sWzAsNSwiIiwxLHsib2Zmc2V0IjoxLCJzdHlsZSI6eyJoZWFkIjp7Im5hbWUiOiJub25lIn19fV1d
    % added shift left and right to the arrows from 1 to 1
    \[\begin{tikzcd}
        1 & N & N \\
        1 & A & A
        \arrow["0", from=1-1, to=1-2]
        \arrow[shift left=0.20ex, -{}, from=1-1, to=2-1]
        \arrow[shift right=0.20ex, -{}, from=1-1, to=2-1]
        \arrow["s", from=1-2, to=1-3]
        \arrow["h"', dashed, from=1-2, to=2-2]
        \arrow["h"', dashed, from=1-3, to=2-3]
        \arrow["a", from=2-1, to=2-2]
        \arrow["f", from=2-2, to=2-3]
    \end{tikzcd}.\]
    i.e., $h \circ 0 = a$ and $h \circ s = f \circ h$.
\end{definition}

If we require only the existence of an arrow that makes the diagram above commute, we obtain the following.

\begin{definition}[Weak natural number object]
    Given a cartesian closed category $\mathcal{C}$, a \emph{weak natural number object} (WNNO) is an object $N$ together with two arrows $1 \xrightarrow{0} N$ and $N \xrightarrow{s} N$ in $\mathcal C$ such that,
    for all diagrams of the form $1 \xrightarrow{a} A \xrightarrow{b} A$ in $\mathcal{C}$, there exists an arrow $h : N \to A$ that makes the following diagram commute:
    \[\begin{tikzcd}
        1 & N & N \\
        1 & A & A
        \arrow["0", from=1-1, to=1-2]
        \arrow[shift left=0.20ex, -{}, from=1-1, to=2-1]
        \arrow[shift right=0.20ex, -{}, from=1-1, to=2-1]
        \arrow["s", from=1-2, to=1-3]
        \arrow["h"', dashed, from=1-2, to=2-2]
        \arrow["h"', dashed, from=1-3, to=2-3]
        \arrow["a", from=2-1, to=2-2]
        \arrow["f", from=2-2, to=2-3]
    \end{tikzcd}.\]
\end{definition}

\begin{example}[NNO in $\textsf{Set}$]
    The category of sets $\textsf{Set}$, whose terminal object is any singleton set, has a natural number object given by the set of natural numbers $\mathbb{N}$ together with the arrows $0 : 1 \to \mathbb{N}$ and $s : \mathbb{N} \to \mathbb{N}$
    defined by $0(*) = 0$ and $s(n) = n + 1$, respectively. The existence and uniqueness of the arrow $h : \mathbb{N} \to A$ is guaranteed by the principle of countable recursion:
    \[ h(n) := \begin{cases}
        h(0) = a \\
        h(s(n)) = f(h(n))
    \end{cases}. \]
\end{example}

\begin{example}[WNNO which is not an NNO]
    Let $\textsf{Pos}$ be the category of partially ordered sets with monotone maps. Its terminal object, up to isomorphism, is the one-element poset, its products are cartesian products equipped with the lexicographic order,
    and its exponential objects are sets of monotone maps equipped equipped with the pointwise order:
    \[  f \leq_{B^A} g \quad\overset{\textnormal{def}}{\Longleftrightarrow}\quad f(a) \leq_B g(a) \quad\text{for every } a \in A. \]
    The usual evaluation map is monotone, idem the currying map and (CCC1) + (CCC2) are easily verified. Thus $\textsf{Pos}$ is cartesian closed.
    Let $\omega_1$ be the first uncountable ordinal and let $W := (\omega_1)_{\mathrm{disc}}$ be its underlying set equipped with the discrete order, rather than with the usual ordinal order; explicitly:
    \[ \alpha \leq_W \beta \quad\Longleftrightarrow\quad \alpha = \beta. \]
    We choose the arrows $0 : 1 \to W : * \mapsto 0$ and $s : W \to W : \alpha \mapsto \alpha + 1$, which is monotone and well-defined since $\omega_1$ is a limit ordinal.\newline
    Let's consider an arbitrary diagram $1 \xrightarrow{a} A \xrightarrow{f} A$ in $\textsf{Pos}$, then we can define $h$ by transfinite recursion as follows:
    \[ h(\alpha) := \begin{cases}
        h(0) = a(*) \\
        h(s(\alpha)) = f(h(\alpha)) \\
        h(\lambda) = a(*) & \text{for $\lambda$ nonzero limit ordinal}
    \end{cases}. \]
    One can easily verify that $h$ is monotone (a map from a discrete-ordered poset it's always monotone). By definition, we have:
    \[ h \circ 0 = a \qquad\text{and}\qquad h \circ s = f \circ h,
    \]
    therefore $(W,0,s)$ is a WNNO in $\textsf{Pos}$.\newline
    This WNNO is not an NNO, in fact, let $2 = \{0,1\}$ have the discrete order, let $a : 1 \to 2$ select $0$, and take $1_2 : 2 \to 2$. Both the constant map $h_1 : W \to 2$ with value $0$ and the map
    \[
        h_2(\alpha) :=
        \begin{cases}
            1 & \text{if } \alpha = \omega + n \text{ for some } n \in \mathbb{N}, \\
            0 & \text{otherwise}
        \end{cases}
    \]
    are monotone and satisfy $h_i \circ 0 = a$ and $h_i \circ s = 1_2 \circ h_i$.
\end{example}

\begin{example}[NNO in $\textsf{Pos}$]
    $(\NN_{\text{disc}}, 0, s)$, where $\NN_{\text{disc}}$ is the set of natural numbers equipped with the discrete order, is an NNO in $\textsf{Pos}$, thanks to the principle of countable recursion.
    The reason why this is unique, whereas it need not to be unique for $(\omega_1)_{\textnormal{disc}}$, is that the commutativity of the diagram fixes $h(0) = a(*)$ and
    recursively defines $h(s(n))$ for every successor of $0 \in N$. Therefore, if we take an ordinal containing a nonzero limit ordinal,
    we obtain an “orbit” under the successor function, determined by a different initial value of $h$; however, the diagram above gives us only two conditions.
\end{example}

\begin{notation}
    In a cartesian closed category $\mathcal{C}$ with a WNNO $(N,0,s)$, we will denote by $J_A(a,f)$ an arrow $h : N \to A$ such that $h \circ 0 = a$ and $h \circ s = f \circ h$.
    Using this notation we can write the following map
    \[ J_A : \Hom(1,A) \times \Hom(A,A) \to \Hom(N,A) \]
    which assigns to each pair $(a,f)$ an arrow $J_A(a,f)$ that makes the diagram
    % https://q.uiver.app/#q=WzAsNixbMCwwLCIxIl0sWzAsMSwiMSJdLFsxLDAsIk4iXSxbMiwwLCJOIl0sWzEsMSwiQSJdLFsyLDEsIkEiXSxbMCwxLCIiLDAseyJvZmZzZXQiOjEsInN0eWxlIjp7ImhlYWQiOnsibmFtZSI6Im5vbmUifX19XSxbMSwwLCIiLDAseyJvZmZzZXQiOjEsInN0eWxlIjp7ImhlYWQiOnsibmFtZSI6Im5vbmUifX19XSxbMCwyLCIwIl0sWzIsMywicyJdLFsxLDQsImEiXSxbNCw1LCJmIl0sWzIsNCwiSl9BKGEsZikiLDIseyJzdHlsZSI6eyJib2R5Ijp7Im5hbWUiOiJkYXNoZWQifX19XSxbMyw1LCJKX0EoYSxmKSIsMix7InN0eWxlIjp7ImJvZHkiOnsibmFtZSI6ImRhc2hlZCJ9fX1dXQ==
    \[\begin{tikzcd}
        1 & N & N \\
        1 & A & A
        \arrow["0", from=1-1, to=1-2]
        \arrow[shift left=0.20ex, -{}, from=1-1, to=2-1]
        \arrow[shift right=0.20ex, -{}, from=1-1, to=2-1]
        \arrow["s", from=1-2, to=1-3]
        \arrow["{J_A(a,f)}"', dashed, from=1-2, to=2-2]
        \arrow["{J_A(a,f)}"', dashed, from=1-3, to=2-3]
        \arrow["a", from=2-1, to=2-2]
        \arrow["f", from=2-2, to=2-3]
    \end{tikzcd}\]
    commute, i.e.,\ $J_A(a,f) \circ 0 = a$ and $J_A(a,f) \circ s = f \circ J_A(a,f)$. [DISCUTERE BREVEMENTE QUANTO AC SERVA]
\end{notation}

Let's see now that if a cartesian closed category has a NNO/WNNO, then the polynomial category over it has the same NNO/WNNO.

\begin{proposition}[NNO/WNNO in polynomial categories]
    If a cartesian closed category $\mathcal{C}$ has a natural number object (resp. a weak natural number object), and if $x : 1 \to A$
    is an indeterminate arrow, then the cartesian closed category $\mathcal{C}[x]$ has the same natural number object (resp. weak natural number object).
\end{proposition}

\begin{proof}
    Let $(N,0,s)$ be a NNO (resp. WNNO) in $\mathcal{C}$, and let's consider the following diagram in $\mathcal{C}[x]$:
    \[ 1 \xrightarrow{\beta(x)} B \xrightarrow{\varphi(x)} B. \]
    We're looking for a polynomial arrow $\chi(x) : N \to B$ such that the following diagram commutes in $\mathcal{C}[x]$:
    % https://q.uiver.app/#q=WzAsNixbMCwwLCIxIl0sWzEsMCwiTiJdLFsyLDAsIk4iXSxbMCwxLCIxIl0sWzEsMSwiQiJdLFsyLDEsIkIiXSxbMCwxLCIwIl0sWzEsMiwicyJdLFs0LDUsIlxcdmFycGhpKHgpIiwyXSxbMSw0LCJcXGNoaSh4KSIsMix7InN0eWxlIjp7ImJvZHkiOnsibmFtZSI6ImRhc2hlZCJ9fX1dLFsyLDUsIlxcY2hpKHgpIiwyLHsic3R5bGUiOnsiYm9keSI6eyJuYW1lIjoiZGFzaGVkIn19fV0sWzMsNCwiXFxiZXRhKHgpIiwyXSxbMCwzLCIiLDIseyJvZmZzZXQiOjEsInN0eWxlIjp7ImhlYWQiOnsibmFtZSI6Im5vbmUifX19XSxbMCwzLCIiLDEseyJvZmZzZXQiOi0xLCJzdHlsZSI6eyJoZWFkIjp7Im5hbWUiOiJub25lIn19fV1d
    % shift left and right to the arrows from 1 to 1
    \[\begin{tikzcd}
        1 & N & N \\
        1 & B & B
        \arrow["0", from=1-1, to=1-2]
        \arrow[shift right = 0.20ex, no head, from=1-1, to=2-1]
        \arrow[shift left = 0.20ex, no head, from=1-1, to=2-1]
        \arrow["s", from=1-2, to=1-3]
        \arrow["{\chi(x)}"', dashed, from=1-2, to=2-2]
        \arrow["{\chi(x)}"', dashed, from=1-3, to=2-3]
        \arrow["{\beta(x)}"', from=2-1, to=2-2]
        \arrow["{\varphi(x)}"', from=2-2, to=2-3]
    \end{tikzcd},\]
    i.e.,\, $\chi(x) \circ 0 \underset{x}{=} \beta(x)$ and $\chi(x) \circ s \underset{x}{=} \varphi(x) \circ \chi(x)$. Applying the functional completeness theorem [ADD REFERENCE] to these polynomial arrow we obtain:
    \begin{itemize}
        \item $b := \kappa_{x \in A}(\beta(x)) : A \times 1 \to B$;
        \item $f := \kappa_{x \in A}(\varphi(x)) : A \times B \to B$;
        \item $h := \kappa_{x \in A}(\chi(x)) : A \times N \to B$.
    \end{itemize}
    Now, by the uniqueness given by functional completeness, finding $\chi(x)$ is equivalent to finding $h$. We can also apply the functional completeness
    to the commutativity of the diagram above, obtaining:
    \[ h \circ \langle \pi_{A,1}, 0 \circ \pi_{A,1}' \rangle = b \qquad \text{and} \qquad h \circ \langle \pi_{A,N}, s \circ \pi_{A,N}' \rangle = f \circ \langle \pi_{A,N}, h \rangle, \]
    or, equivalently:
    \[ h \circ (1_A \times 0) = b \qquad \text{and} \qquad h \circ (1_A \times s) = f \circ \langle \pi_{A,N}, h \rangle. \]
    Let $\sigma_{1,A} := \langle \pi_{1,A}', \pi_{1,A} \rangle : 1 \times A \to A \times 1$ be the isomorphism given by the symmetry of the product, we can take the first equation, compose it with $\sigma_{1,A}$ and apply the currying isomorphism to obtain:
    \[  \Phi(h \circ (1_A \times 0) \circ \sigma_{1,A}) = \Phi(b \circ \sigma_{1,A}) =: b' : 1 \to B^A, \]
    since both $\Phi$ is a bijection, $b'$ is uniquely determined by $b$. In the LHS, we have $(1_A \times 0) \circ \sigma_{1,A} = \sigma_{N,A} \circ (0 \times 1_A)$, where $\sigma_{N,A} : N \times A \to A \times N$ is the symmetry of the product.
    Moreover, by the distributivity of the currying, $\Phi(h \circ \sigma_{N,A} \circ (0 \times 1_A)) = \Phi(h \circ \sigma_{N,A}) \circ 0$, therefore if we write $h' := \Phi(h \circ \sigma_{N,A}) : N \to B^A$, which is, as before, uniquely determined by $h$, we find that the first equation becomes:
    \[ h' \circ 0 = b'. \]
    Similarly, we precompose the second equation with $\sigma_{N,A} : N \times A \to A \times N$ and apply the currying, this way we get arrows $N \to B^A$ on both sides. The LHS becomes
    \[ \Phi(h \circ (1_A \times s) \circ \sigma_{N,A}) = \Phi(h \circ \sigma_{N,A} \circ (s \times 1_A)) = h' \circ s, \]
    where $(1_A \times s) \circ \sigma_{N,A} = \sigma_{N,A} \circ (s \times 1_A)$ follows form the definition of $\sigma_{N,A}$. The RHS becomes
    \[ \Phi(f \circ \langle \pi_{A,N}, h \rangle \circ \sigma_{N,A}), \]
    heuristically, we would like to get $f' \circ h'$, for some definition of $f' : B^A \to B^A$.
    If we find a map $q$ such that $q \circ (h' \times 1_A) = f \circ \langle \pi_{A,N}, h \rangle \circ \sigma_{N,A}$, then we can use the distributivity of the currying to get $\Phi(q) \circ h'$.\newline
    Taking $q := f \circ \langle \pi_{B^A,A}', \ev_{B,A}\rangle : B^A \times A \to B$, we have
    \begin{align*}
        q \circ (h' \times 1_A) &= f \circ \langle \pi_{B^A,A}', \ev_{B,A}\rangle \circ (h' \times 1_A) \\
                                &= f \circ \langle \pi_{N,A}',\Psi(h') \rangle \\
                                &= f \circ \langle \pi_{N,A}', \Psi(\Phi(h \circ \sigma_{N,A})) \rangle \\
                                &= f \circ \langle \pi_{N,A}', h \circ \sigma_{N,A} \rangle \\
                                &= f \circ \langle \pi_{A,N}, h \rangle \circ \sigma_{N,A},
    \end{align*}
    therefore, by taking $f' := \Phi(q)$, uniquely determined by $f$, we finally have that the second equation in the form $h' \circ s = f' \circ h'$. We now want the following equations to hold in $\mathcal{C}$:
    \[ h' \circ 0 = b' \qquad\text{and}\qquad h' \circ s = f' \circ h'. \]
    These equations corresponds to the commutativity of the diagram:
    % https://q.uiver.app/#q=WzAsNixbMCwwLCIxIl0sWzEsMCwiTiJdLFsyLDAsIk4iXSxbMSwxLCJCXkEiXSxbMiwxLCJCXkEiXSxbMCwxLCIxIl0sWzAsMSwiMCJdLFsxLDIsInMiXSxbNSwzLCJiJyIsMl0sWzAsNSwiIiwxLHsib2Zmc2V0IjoxLCJzdHlsZSI6eyJoZWFkIjp7Im5hbWUiOiJub25lIn19fV0sWzAsNSwiIiwxLHsib2Zmc2V0IjotMSwic3R5bGUiOnsiaGVhZCI6eyJuYW1lIjoibm9uZSJ9fX1dLFszLDQsImYnIiwyXSxbMSwzLCJoJyIsMix7InN0eWxlIjp7ImJvZHkiOnsibmFtZSI6ImRhc2hlZCJ9fX1dLFsyLDQsImgnIiwwLHsic3R5bGUiOnsiYm9keSI6eyJuYW1lIjoiZGFzaGVkIn19fV1d
    % added shift left and right to the arrows from 1 to 1
    \[\begin{tikzcd}
        1 & N & N \\
        1 & {B^A} & {B^A}
        \arrow["0", from=1-1, to=1-2]
        \arrow[shift right = 0.20ex, no head, from=1-1, to=2-1]
        \arrow[shift left = 0.20ex, no head, from=1-1, to=2-1]
        \arrow["s", from=1-2, to=1-3]
        \arrow["{h'}"', dashed, from=1-2, to=2-2]
        \arrow["{h'}", dashed, from=1-3, to=2-3]
        \arrow["{b'}"', from=2-1, to=2-2]
        \arrow["{f'}"', from=2-2, to=2-3]
    \end{tikzcd},\]
    since $(N,0,s)$ is a NNO (resp. WNNO) in $\mathcal{C}$, there exists an arrow $h' : N \to B^A$ that makes the diagram above commute, and by the functional completeness and the bijection used we can uniquely determine $\chi(x) : N \to B$ in $\mathcal{C}[x]$ such that the original diagram commutes.\newline
    In the case of a NNO, the uniqueness of $h'$ implies the uniqueness of $\chi(x)$, again thanks to the functional completeness and the bijections used.
\end{proof}

\begin{remark}[One the choice of $q$]
    The heuristic behind the choice of $q = f \circ \langle \pi_{B^A,A}', \ev_{B,A}\rangle : B^A \times A \to B$ is pointwise. Using for a moment the set-theoretic interpretation, an element $g \in B^A$ is a map $g : A \to B$, and
    \[ q(g,a) = f(a,g(a)). \]
    Therefore, for $(n,a) \in N \times A$, we have
    \[ \big(q \circ (h' \times 1_A)\big)(n,a) = q(h'(n),a) = f(a,h'(n)(a)). \]
    Since $\Psi(h') = h \circ \sigma_{N,A}$, we have $h'(n)(a) = h(a,n)$, that is precisely the value of $f \circ \langle \pi_{A,N}, h \rangle \circ \sigma_{N,A}$ at $(n,a)$.
\end{remark}

We introduced the notation $J_A(a,f)$ to denote a generic arrow that makes the diagram commute, but this notation does not express the
fact that $J_A(a,f)$ is a function which associates to each pair $(a,f)$ an arrow that makes the diagram commute.
Now that we have the proposition above, we can define, for each object, an actual function that associates to each pair $(a,f)$ an arrow that makes the diagram commute, and we can do it in the following way.

\begin{definition}[Parameterized iterator]
    Let $(N,0,s)$ be a WNNO in $\mathcal{C}$, $B$ an object of $\mathcal{C}$ and $x : 1 \to B \times B^B$ an indeterminate arrow.
    By the preceding proposition, $(N,0,s)$ is also a WNNO in $\mathcal{C}[x]$. In this category, let
    \[ y(x) := \pi_{B,B^B} \circ x : 1 \to B \qquad\text{and}\qquad v(x) := \pi_{B,B^B}' \circ x : 1 \to B^B. \]
    Since $v(x)^\sharp : B \to B$, we can choose an arrow $\chi_B(x) : N \to B$ in $\mathcal{C}[x]$ such the following diagram commutes
    % https://q.uiver.app/#q=WzAsNixbMCwwLCIxIl0sWzEsMCwiTiJdLFsyLDAsIk4iXSxbMSwxLCJCIl0sWzIsMSwiQiJdLFswLDEsIjEiXSxbMCwxLCIwIl0sWzEsMiwicyJdLFszLDQsInYoeCleXFxzaGFycCIsMl0sWzAsNSwiIiwyLHsib2Zmc2V0IjoxLCJzdHlsZSI6eyJoZWFkIjp7Im5hbWUiOiJub25lIn19fV0sWzAsNSwiIiwyLHsib2Zmc2V0IjotMSwic3R5bGUiOnsiaGVhZCI6eyJuYW1lIjoibm9uZSJ9fX1dLFs1LDMsInkoeCkiLDJdLFsxLDMsIlxcY2hpKHgpIiwyLHsic3R5bGUiOnsiYm9keSI6eyJuYW1lIjoiZGFzaGVkIn19fV0sWzIsNCwiXFxjaGkoeCkiLDIseyJzdHlsZSI6eyJib2R5Ijp7Im5hbWUiOiJkYXNoZWQifX19XV0=
    % shift left and right to the arrows from 1 to 1
    \[\begin{tikzcd}
        1 & N & N \\
        1 & B & B
        \arrow["0", from=1-1, to=1-2]
        \arrow[shift right = 0.20ex, no head, from=1-1, to=2-1]
        \arrow[shift left = 0.20ex, no head, from=1-1, to=2-1]
        \arrow["s", from=1-2, to=1-3]
        \arrow["{\chi(x)}"', dashed, from=1-2, to=2-2]
        \arrow["{\chi(x)}"', dashed, from=1-3, to=2-3]
        \arrow["{y(x)}"', from=2-1, to=2-2]
        \arrow["{v(x)^\sharp}"', from=2-2, to=2-3]
    \end{tikzcd},\]
    i.e.,
    \[ \chi_B(x) \circ 0 \underset{x}{=} y(x) \qquad\text{and}\qquad
    \chi_B(x) \circ s \underset{x}{=} v(x)^\sharp \circ \chi_B(x). \]
    Functional completeness associates with this choice the unique arrow
    \[ I_B := \kappa_{x \in B \times B^B}(\chi_B(x)) : (B \times B^B) \times N \to B \]
    in $\mathcal{C}$ such that
    \[ I_B \circ \langle x \circ !_N,1_N\rangle \underset{x}{=} \chi_B(x), \]
    we call $I_B$ the \emph{parameterized iterator} for $B$.
\end{definition}

\begin{remark}
    Following the definition, if we have $p : 1 \to B \times B^B$ in $\mathcal{C}$, then $h_p := I_B \circ \langle p \circ !_N,1_N\rangle : N \to B$ is an arrow that makes the diagram commute.
    At the moment, it's as if we had a “function" $x \mapsto I_B(x,\cdot)$, which associates to each pair $(b,f)$, encoded by $x$ as $\langle a, \ulcorner f \urcorner \rangle$, an arrow $h : N \to B$ that makes the diagram commute.
    We want to get something a little bit more practical, i.e.,\ a “function" which separately takes $a$ and $f$ as arguments, and returns an arrow $h : N \to B$ that makes the diagram commute.
\end{remark}

Let's consider $y : 1 \to B$, $v : 1 \to B^B$ and $z : 1 \to N$ indeterminate arrows, let's also define $\mathcal{C}[y,v] := \mathcal{C}[y][v]$ and $\mathcal{C}[y,v,z] := \mathcal{C}[y][v][z]$.\newline
The universal property of polynomial categories gives a canonical isomorphism $F : \mathcal{C}[x] \xrightarrow{\sim} \mathcal{C}[y,v]$ over $\mathcal{C}$, determined by $x \mapsto \langle y,v\rangle$;
its inverse sends $y$ to $y(x)$ and $v$ to $v(x)$, therefore $F = \widehat{\iota_{y,v}}$. It's easy to check that the isomorphism extends, in the natural way, to $\mathcal{C}[x][z] \cong \mathcal{C}[y,v,z]$, through $\widehat F$.
Since $F$ is a cartesian closed functor and $\mathcal C[x]$ and $\mathcal C[y,v]$ has the same WNNO, applying $F$ to $\chi_B(x)$ preserves the commutativity of the diagram.
Due to the fact that $v(x) \leftrightarrow v$ and $y(x) \leftrightarrow y$ we have $F(\chi_B(x)) = J_B(y,v^\sharp)$. Using the definition of iterator in $\mathcal{C}[y,v]$, we have
\[ J_B(y,v^\sharp) \underset{y,v}{=} I_B \circ \langle \langle y,v\rangle \circ !_N,1_N\rangle.\]
If we now add the indeterminate $z : 1 \to N$, precomposing with $z$ gives:
\begin{align*}
    J_B(y,v^\sharp) \circ z &\underset{y,v,z}{=} I_B \circ \langle \langle y,v\rangle \circ !_N,1_N\rangle \circ z \\
                            &\underset{y,v,z}{=} I_B \circ \langle \langle y,v\rangle z\rangle.
\end{align*}
Using the notation $\langle y,v,z\rangle := \langle \langle y,v\rangle, z\rangle : 1 \to (B \times B^B) \times N$ we get the equality in $\mathcal{C}[y,v,z]$:
\[ J_B(y,v^\sharp) \circ z \underset{y,v,z}{=} I_B \circ \langle y,v,z\rangle. \]

\begin{corollary}[Internal form of a WNNO]
    An object $N$ equipped with arrows
    \[ 1 \xrightarrow{0} N \xrightarrow{s} N \]
    is a WNNO in a cartesian closed category $\mathcal{C}$ if and only if, for every object $B$ of $\mathcal{C}$, there exists an arrow $I_B : (B \times B^B) \times N \to B$ such that, with the preceding notation for $y$, $v$ and $z$,
    \begin{align*}
        I_B \circ \langle y,v,0\rangle &\underset{y,v}{=} y, \\
        I_B \circ \langle y,v,s \circ z\rangle &\underset{y,v,z}{=} v^\sharp \circ I_B \circ \langle y,v,z\rangle.
    \end{align*}
\end{corollary}

\begin{proof}
    Let's prove the two implications separately.
    \begin{itemize}
        \item[$\boxed{\Longrightarrow}$] Let's assume that $(N,0,s)$ is a WNNO in $\mathcal{C}$. By the construction above, we have the identity
        \[ J_B(y,v^\sharp) \circ z \underset{y,v,z}{=} I_B \circ \langle y,v,z\rangle, \]
        in $\mathcal{C}[y,v,z]$. If we apply the substitution functor $S[0/z] : \mathcal{C}[y,v,z] \to \mathcal{C}[y,v]$ to both sides, we get
        \[ J_B(y,v^\sharp) \circ 0 \underset{y,v}{=} I_B \circ \langle y,v,0\rangle. \]
        Since $J_B(y,v^\sharp) \circ 0 \underset{y,v}{=} y$, we have
        \[ y \underset{y,v}{=} I_B \circ \langle y,v,0\rangle. \]
        Similarly, let's temporarily use $t : 1 \to N$ instead of $z$ in the identity above, and apply the substitution functor $S[s \circ z/t] : \mathcal{C}[y,v,t] \to \mathcal{C}[y,v,z]$ to both sides, we get
        \[ J_B(y,v^\sharp) \circ s \circ z \underset{y,v,z}{=} I_B \circ \langle y,v,s \circ z\rangle. \]
        Using the fact that $J_B(y,v^\sharp) \circ s \underset{y,v}{=} v^\sharp \circ J_B(y,v^\sharp)$, we have
        \[ v^\sharp \circ J_B(y,v^\sharp) \circ z \underset{y,v,z}{=} I_B \circ \langle y,v,s \circ z\rangle, \]
        and, finally, using again the initial identity, $J_B(y,v^\sharp) \circ z \underset{y,v,z}{=} I_B \circ \langle y,v,z\rangle$, we find
        \[ v^\sharp \circ I_B \circ \langle y,v,z\rangle \underset{y,v,z}{=} I_B \circ \langle y,v,s \circ z\rangle. \]
        In both cases we used the commutativity of the diagram given by the initial WNNO in $\mathcal C[y,v]$, which also holds in $\mathcal C[y,v,z]$.
        \item[$\boxed{\Longleftarrow}$] Conversely, suppose that for every object $B$ of $\mathcal{C}$ there exists an arrow $I_B : (B \times B^B) \times N \to B$ such that the two polynomial identities in the statement are true.
        Given a diagram
        % https://q.uiver.app/#q=WzAsNixbMCwxLCIxIl0sWzEsMSwiQiJdLFsyLDEsIkIiXSxbMCwwLCIxIl0sWzEsMCwiTiJdLFsyLDAsIk4iXSxbMyw0LCIwIl0sWzQsNSwicyJdLFswLDEsImIiLDJdLFsxLDIsImYiLDJdLFszLDAsIiIsMix7Im9mZnNldCI6MSwic3R5bGUiOnsiaGVhZCI6eyJuYW1lIjoibm9uZSJ9fX1dLFszLDAsIiIsMSx7Im9mZnNldCI6LTEsInN0eWxlIjp7ImhlYWQiOnsibmFtZSI6Im5vbmUifX19XV0=
        % shift left and right to the arrows from 1 to 1
        \[\begin{tikzcd}
            1 & N & N \\
            1 & B & B
            \arrow["0", from=1-1, to=1-2]
            \arrow[shift right = 0.20ex, no head, from=1-1, to=2-1]
            \arrow[shift left = 0.20ex, no head, from=1-1, to=2-1]
            \arrow["s", from=1-2, to=1-3]
            \arrow["b"', from=2-1, to=2-2]
            \arrow["f"', from=2-2, to=2-3]
        \end{tikzcd},\]
        we need to find an arrow $h : N \to B$ that makes it commute. Let $\langle b,\ulcorner f\urcorner\rangle : 1 \to B \times B^B$, then we have the arrow
        \[ I_B \circ \langle \langle b,\ulcorner f\urcorner\rangle \circ !_N,1_N\rangle = I_B \circ \langle b \circ !_N, \ulcorner f\urcorner \circ !_N, 1_N\rangle : N \to B. \]
        Now, precomposing and using the two polynomial identities we get
        \begin{align*}
            I_B \circ \langle b \circ !_N, \ulcorner f\urcorner \circ !_N, 1_N\rangle \circ 0 &= I_B \circ \langle b, \ulcorner f\urcorner, 0\rangle \\
                                                                                              &= b,
        \end{align*}
        and, similarly
        \begin{align*}
            I_B \circ \langle b \circ !_N, \ulcorner f\urcorner \circ !_N, 1_N\rangle \circ s \circ z &\underset{z}{=} I_B \circ \langle b, \ulcorner f\urcorner, s \circ z\rangle \\
                                                                                                      &\underset{z}{=} \ulcorner f\urcorner^\sharp \circ I_B \circ \langle b, \ulcorner f\urcorner, z\rangle \\
                                                                                                      &\underset{z}{=} f \circ I_B \circ \langle b, \ulcorner f\urcorner, z\rangle \\
                                                                                                      &\underset{z}{=} f \circ I_B \circ \langle b \circ !_N, \ulcorner f\urcorner \circ !_N, 1_N\rangle \circ z.
        \end{align*}
        Therefore, $h := I_B \circ \langle b \circ !_N, \ulcorner f\urcorner \circ !_N, 1_N\rangle$ makes the diagram commute pointwise
    \end{itemize}
\end{proof}
%DOVREI AGGIUNGERE LA DEFINIZIONE DI FUNTORE CARTESIANO CHIUSO CHE PRESERVA NNO/WNNO(?)

\section{Typed \texorpdfstring{$\lambda$}{lambda}-calculi}
Now that we have seen cartesian closed categories with WNNOs, we know categories that are models of positive intuitionistic propositional
logic, and have a notion of recursion. We also want to see how this kind of categories can interpret typed $\lambda$-calculus,
and to see the latter as a sort of “internal language” of cartesian closed categories.

The specific kind of typed $\lambda$-calculus we will consider is a \emph{typed $\lambda\beta\eta$-calculus}, with \emph{terminal type}, \emph{product types}, \emph{surjective pairing} and \emph{iterator}.
This will be the first step towards establishing Lambek's equivalence of categories.

\begin{definition}
    A \emph{typed $\lambda$-calculus} is a triple $(\textsf{Ty},\textsf{Tm},\textsf{Eq})$, made of \emph{types}, \emph{terms} and \emph{equations} between terms,
    which are said to \emph{hold}, all subject to certain closure conditions. We write $a \in A$ to say that $a$ is a term of type $A$.
    The types are defined recursively as follows:
    \begin{itemize}
        \item[(Ty1)] there are $1,N$ in $\textsf{Ty}$, called \emph{basic types};
        \item[(Ty2)] if $A$ and $B$ are types, then $A \times B$ is a type;
        \item[(Ty3)] if $A$ and $B$ are types, then $B^A$ is a type.
    \end{itemize}
    The class of terms is defined recursively as follows:
    \begin{itemize}
        \item[(Tm1)] for each type $A$, there are countably many variables, $x_i^A \in A$, of type $A$;
        \item[(Tm2)] the type $1$ has a term $* \in 1$;
        \item[(Tm3a)] if $a \in A$ and $b \in B$, then there is a term $\langle a,b\rangle$ of type $A \times B$;
        \item[(Tm3b)] if $c \in A \times B$, then there are $\pi_{A,B}(c) \in A$ and $\pi_{A,B}'(c) \in B$;
        \item[(Tm4)] if $f \in B^A$ and $a \in A$, then $\ev_{A,B}(f,a)$ is a term of type $B$;
        \item[(Tm5)] if $x \in A$ is a variable and $\varphi(x) \in B$ is a term, then $\lambda_{x \in A} \varphi(x) \in B^A$;\newline (\emph{$\lambda$-abstraction})
        \item[(Tm6a)] there is a term $0 \in N$;
        \item[(Tm6b)] if $n \in N$, then $s(n) \in N$;
        \item[(Tm7)] if $a \in A$, $f \in A^A$ and $n \in N$, then $I_A(a,f,n) \in A$.
    \end{itemize}
    In (Tm5) we mean that there can be a free occurrence of $x$ in $\varphi(x)$.
    We'll often abbreviate $\ev_{A,B}(f,a)$ as $f^\sharp(a)$, when the types are clear from the context.\\
    We now define recursively the class of equations between terms:
    \begin{itemize}
        \item[(Eq0)] every equation between terms is of the form $a \underset{X}{=} b$, where $a$ and $b$ are terms of the same type,
        and $X$ is a finite set of variables containing all the free variables of $a$ and $b$;
        \item[(Eq1)] $\underset{X}{=}$ is a binary relation on terms of the same type, specifically, it is an equivalence relation;
        \item[(Eq2)] if $a \underset{X}{=} b$ and $X \subseteq Y$, then $a \underset{Y}{=} b$;
        \item[(Eq3a)] if $a \underset{X}{=} b$, then $f^\sharp(a) \underset{X}{=} f^\sharp(b)$, for any $f \in B^A$ and $a,b \in A$;
        \item[(Eq3b)] if $\varphi(x) \underset{X \cup\{x\}}{=} \varphi'(x)$, then $\lambda_{x \in A} \varphi(x) \underset{X}{=} \lambda_{x \in A} \varphi'(x)$, for a variable $x \in A$ and $\varphi(x),\varphi'(x) \in B$;
        \item[(Eq4)] $a \underset{X}{=} *$, for any $a \in 1$;
        \item[(Eq5a)] $\pi_{A,B}(\langle a,b\rangle) \underset{X}{=} a$, for all $a \in A$ and $b \in B$;
        \item[(Eq5b)] $\pi_{A,B}'(\langle a,b\rangle) \underset{X}{=} b$, for all $a \in A$ and $b \in B$;
        \item[(Eq5c)] $\langle \pi_{A,B}(c),\pi_{A,B}'(c)\rangle \underset{X}{=} c$, for all $c \in A \times B$; (\emph{surjective pairing})
        \item[(Eq6a)] $(\lambda_{x \in A} \varphi(x))^\sharp(a) \underset{X}{=} \varphi(a)$, for all $a \in A$, which are substitutable for $x$ in $\varphi(x)$; (\emph{$\beta$-reduction})
        \item[(Eq6b)] $\lambda_{x \in A}(f^\sharp(x)) \underset{X}{=} f$, for all $f \in B^A$ such that $x \notin X$; (\emph{$\eta$-reduction})
        \item[(Eq7a)] $I_A(a,f,0) \underset{X}{=} a$, for all $a \in A$, $f \in A^A$ and $0 \in N$;
        \item[(Eq7b)] $I_A(a,f,s(n)) \underset{X}{=} f^\sharp(I_A(a,f,n))$, for all $a \in A$, $f \in A^A$, $n \in N$ and $s(n) \in N$;
        \item[(Eq8)] $\lambda_{x \in A} \varphi(x) \underset{X}{=} \lambda_{x' \in A} \varphi(x')$, if $x'$ is substitutable for $x$ in $\varphi(x)$ and $x'$ is not free in $\varphi(x)$; (\emph{$\alpha$-conversion}).
    \end{itemize}
\end{definition}

\begin{remark}[On the definition of typed $\lambda$-calculi]
    There may be other basic types, terms, type-forming rules, term-forming rules and equations.
    What we defined is a \emph{simply typed $\lambda$-calculus} with product types and iterator, that any other typed $\lambda$-calculus, also with product types and iterator, can be seen as an extension of.
\end{remark}

\begin{remark}
    (Eq7), the $\alpha$-conversion rule, may be omitted if we are willing to identify terms that differ only by the names of their bound variables.
\end{remark}

\begin{notation}[Free variables and closed terms]
    We denote the free variables of a term $t$, with $\FV(t)$. The definition of free variables is the standard one, see for example \cite{Carpenter1997}.
    We also say that a term $t$ is \emph{closed} if $\FV(t) = \varnothing$, and that a term $t$ is \emph{substitutable} for a variable $x$ in another term $\varphi(x)$,
    if no free occurrence of a variable in $t$ becomes bound in $\varphi(t)$.
\end{notation}

The rule (Eq2) already allows us to pass from $a \underset{X}{=} b$ to $b \underset{X \cup \{x\}}{=} a$.
Under certain assumptions, we can go the opposite way, as we see in the following propositions.

\begin{proposition}[Equality implies pointwise equality]
    In any typed $\lambda$-calculus, if $\varphi(x) \underset{X \cup \{x\}}{=} \psi(x)$, for $x \in A$, then $\varphi(a) \underset{X}{=} \psi(a)$, for all $a \in A$, assuming that:
    \begin{enumerate}[label=(\roman*)]
        \item $x \not \in X$;
        \item $\FV(a) \subseteq X$, for all $a \in A$, otherwise the thesis does not make sense.
    \end{enumerate}
\end{proposition}

\begin{proof}
    From the hypothesis, we apply the rule (Eq3b) to get
    \[ \lambda_{x \in A} \varphi(x) \underset{X}{=} \lambda_{x \in A} \psi(x), \]
    and then, applying the $\beta$-reduction rule to both sides, we get
    \[ \varphi(a) \underset{X}{=} (\lambda_{x \in A} \varphi(x))^\sharp(a) \underset{X}{=} (\lambda_{x \in A} \psi(x))^\sharp(a) \underset{X}{=} \psi(a). \]
\end{proof}

\begin{corollary}
    If $f$ and $g$ do not contain free occurrences of $x \in A$, then $f(x) \underset{X \cup \{x\}}{=} g(x)$ implies $f \underset{X}{=} g$,
    assuming that the type is \emph{inhabited}, i.e., there exists a closed term $a \in A$.
\end{corollary}

\begin{proof}
   Let $a \in A$ be a closed term, then from $f(x) \underset{X \cup \{x\}}{=} g(x)$, using the previous proposition, we get $f(a) \underset{X}{=} g(a)$.
   Since no free occurrence of $x$ is in $f$ and $g$, we have $f(a) \underset{X}{=} f$ and $g(a) \underset{X}{=} g$ (formally by using the definition of substitution), therefore, $f \underset{X}{=} g$.
\end{proof}

It can happen that the type $A$ is \emph{empty}, or, on the other hand, that every type is inhabited, and, therefore, the previous corollary is always applicable.\newline
If we know that $f$ and $g$ are of type $B^A$, then we prove the corollary above, without the assumption that $A$ is inhabited, as follows.

\begin{proposition}
    Given $f,g \in B^A$, if $f \underset{X \cup \{x\}}{=} g$, for $x \in A$, then $f \underset{X}{=} g$.
\end{proposition}

\begin{proof}
    We'd like to use $\eta$-reduction, but we first need to go from the hypothesis to $f^\sharp(x) \underset{X \cup \{x\}}{=} g^\sharp(x)$.
    In order to do that, let's first define a new term. Let $h \in B^A$ and $x \in A$ a variable, then $h^\sharp(x) \in B$, and we can define the term:
    \[ \varphi(x) = \lambda_{h \in B^A} h^\sharp(x) \in B^{B^A}. \]
    Since $f \underset{X \cup \{x\}}{=} g$, we can apply (Eq3b) to get
    \[ \varphi^\sharp(f) \underset{X \cup \{x\}}{=} \varphi^\sharp(g), \]
    and then, by $\beta$-reduction, we have
    \[ f^\sharp(x) \underset{X \cup \{x\}}{=} g^\sharp(x). \]
    Now we can apply the $\lambda$-abstraction to both sides, since they are $A$-terms, hence
    \[ \lambda_{x \in A} (f^\sharp(x)) \underset{X}{=} \lambda_{x \in A} (g^\sharp(x)). \]
    Finally, by $\eta$-reduction, we get $f \underset{X}{=} g$.
\end{proof}

Let's see some examples of typed $\lambda$-calculi.

\begin{example}[Simply typed $\lambda$-calculus]
    As we already mentioned, if there are not types, terms, equations, type-forming rules and term-forming rules other than those in the definition above,
    we have what is called a \emph{simply typed $\lambda\beta\eta$-calculus} with terminal type, product types, surjective pairing and iterator.
    Also called \emph{pure typed $\lambda$-calculus} with \emph{weak natural object} and all the other features above. We denote this calculus by $\mathcal L_0$.
\end{example}

\begin{example}[$\lambda$-calculus generated by a graph]

\end{example}

The next example is fundamental for this thesis, and it's perfect as a bridge for the next section.

\begin{example}[Internal language of a CCC]

\end{example}

\section{Lambek's equivalence}
